% TOP_PROBLEM: {"id": 30006791, "title": "Brauer-Manin obstruction for number-field curves: Hasse-principle form", "category_id": 1, "category": {"id": 1, "name": "number_theory", "display_name": "Number Theory", "description": "Properties of integers, prime numbers, Diophantine equations.", "slug": "number-theory", "order_index": 1, "created_at": "2026-07-31T15:26:25.670Z"}, "status": "open"}
% FIRST_POSTED: 2026-09-27
% Standalone research entry 210; batch 208--217.
% Original section material preserved from Pasted text(3).txt.
% Research additions: 27 September 2026.
% Compile twice with: lualatex 30006791.tex
% !TeX program = lualatex
% Compile twice: lualatex open_problems_500.tex
% All references and prose are embedded; no BibTeX database or external inputs.
\documentclass[11pt,a4paper]{article}
\usepackage[margin=24mm,headheight=14pt]{geometry}
\usepackage{fontspec}
\setmainfont{Latin Modern Roman}
\setsansfont{Latin Modern Sans}
\setmonofont{Latin Modern Mono}
\usepackage{amsmath,amssymb,mathtools}
\usepackage{unicode-math}
\setmathfont{Latin Modern Math}
\usepackage{microtype}
\usepackage{enumitem,xurl,needspace}
\usepackage[dvipsnames]{xcolor}
\usepackage{hyperref}
\hypersetup{unicode=true,colorlinks=true,linkcolor=MidnightBlue,urlcolor=MidnightBlue,
 pdftitle={500 Mathematical Problem Statements},pdfauthor={Source-annotated compilation},bookmarksnumbered=true}
\usepackage{bookmark}
\usepackage{tocloft}
\setlength{\cftsecnumwidth}{3em}
\cftsetpnumwidth{2.5em}
\cftsetrmarg{4em}
\usepackage{titlesec}
\titleformat{\section}{\Large\bfseries\raggedright}{\thesection}{0.7em}{}
\usepackage{fancyhdr}
\pagestyle{fancy}\fancyhf{}
\fancyhead[L]{\small\sffamily Mathematical problem statements}
\fancyhead[R]{\small Source edition 22}
\fancyfoot[C]{\thepage}
\setlength{\parindent}{0pt}
\setlength{\parskip}{5pt plus 1pt}
\setlength{\emergencystretch}{3em}
\setcounter{tocdepth}{1}
\setcounter{secnumdepth}{1}
\setlist[itemize]{leftmargin=3em,itemsep=5pt,topsep=4pt,parsep=0pt}
\allowdisplaybreaks
\newcommand{\N}{\mathbb{N}}
\newcommand{\Z}{\mathbb{Z}}
\newcommand{\Q}{\mathbb{Q}}
\newcommand{\R}{\mathbb{R}}
\newcommand{\C}{\mathbb{C}}
\newcommand{\F}{\mathbb{F}}
\newcommand{\A}{\mathbb{A}}
\newcommand{\PP}{\mathbb P}
\newcommand{\E}{\mathbb{E}}
\newcommand{\eps}{\varepsilon}
\newcommand{\dd}{\,\mathrm d}
\newcommand{\sref}[2]{\hyperlink{source:#1:#2}{[S#2]}}
\newcommand{\eref}[2]{\hyperlink{extra:#1:#2}{[E#2]}}
% Turn the next switch off for a shorter reading copy. The quotations preserve
% every exactTarget field, including qualifications omitted from a short summary.
\newif\ifshowcatalogue
\showcataloguetrue
\newcommand{\cataloguescope}[1]{\ifshowcatalogue
 \paragraph{Exact scope in the supplied catalogue.}
 \begin{quote}\small #1\end{quote}\fi}

% Standalone wrapper; no external inputs or bibliography files are required.
\fancyhead[L]{\small\sffamily Research notebook: section 30006791}
\fancyhead[R]{\small 27 September 2026}
\hypersetup{pdftitle={Research attempt 210},pdfauthor={Research notebook}}
\setlength{\parskip}{4pt plus 1pt}
\setlist[itemize]{before=\small\interlinepenalty10000,itemsep=3pt}
\begin{document}
\setcounter{section}{0}
% ================================================================
% problemId: 30006791
% Canonical problem ID: 30006791
% exactTarget SHA-256: 6840b7480166c57af81293c20327083d0780c667c9e65d128523823db53225ae
\clearpage
\section*{\texorpdfstring{Brauer-Manin obstruction for number-field curves: Hasse-principle form}{Brauer-Manin obstruction for number-field curves: Hasse-principle form}}\label{30006791}
\subsection{Definitions and mathematical statement}
Let $C/k$ be a smooth projective geometrically connected curve of genus at least two over a number field. Its Brauer--Manin set is
\[
 C(\A_k)^{\operatorname{Br}}=
 \{(P_v):\sum_v\operatorname{inv}_v(\alpha(P_v))=0
 \text{ for every }\alpha\in H^2_{\mathrm{et}}(C,\mathbf G_m)\}.
\]
The local invariants take values in $\Q/\Z$ and the sum has only finitely many nonzero terms for fixed $\alpha,(P_v)$. The selected assertion is nonemptiness of this set implies $C(k)\ne\varnothing$. This is only the Hasse-principle form; density of rational points in the adelic set is not added. No rational divisor class of degree one or finiteness assumption on a Tate--Shafarevich group is permitted as an unstated premise.
\par\smallskip\textit{Formulation and concept references:} \sref{30006791}{1}.
\cataloguescope{For every number field k and every smooth projective geometrically connected curve C/k of genus at least two, if C(A\_k)\textasciicircum{}Br is nonempty then C(k) is nonempty. Here Br(C)=H\textasciicircum{}2\_et(C,G\_m) is the full cohomological Brauer group, and C(A\_k)\textasciicircum{}Br consists of adelic tuples (P\_v) for which sum\_v inv\_v(alpha(P\_v))=0 for every alpha in Br(C). Constants obey global reciprocity. Passing at archimedean places to connected components does not change nonemptiness. No rational base point, degree-one divisor class, Sha finiteness, rank bound or effective Jacobian presentation is assumed.}
\subsection{Short English statement}
For a number-field curve of genus at least two, do adelic points passing every Brauer–Manin test guarantee an actual rational point?
% BEGIN RESEARCH ADDITION 210
\subsection{Research attempt}

\paragraph{Current frontier.}
Stoll identifies the full Brauer set of a curve with its finite abelian
descent set. His Corollary 7.4, assuming the Jacobian's Tate--Shafarevich
group has no nonzero divisible subgroup, both treats the absence of a
degree-one divisor class and describes the Brauer set using the closure
of the Mordell--Weil group. His main theorem also proves the assertion for
curves admitting suitable nonconstant maps to rank-zero abelian varieties.
\sref{30006791}{1} Creutz's September 2026 revision proves an adelic
Mordell--Lang statement over global \emph{function fields} with a
nonisotriviality hypothesis; its analogous number-field results retain
additional conjectural input. It is not the unrestricted number-field
theorem requested here.\eref{30006791}{1}

\paragraph{Attack route.}
The tempting compactness argument produces a profinite Mordell--Weil point,
not an actual rational point. I try to repair it by coupling the finite-place
sieve to an archimedean height budget. Two precise mechanisms are examined:
bounded-height representatives of increasingly many congruences, and a
product-formula inequality forcing a sufficiently small representative
onto the curve. The second yields an unconditional separation lemma over
all number fields; obtaining representatives below its threshold is the
unresolved arithmetic step.

\paragraph{Attempt.}
\textbf{1. Do not assume the missing base point.}
Write $J=\operatorname{Jac}(C)$. Temporarily impose the explicit hypothesis
\begin{equation}\label{30006791:sha}
 \operatorname{Sha}(k,J)_{\mathrm{div}}=0.
\end{equation}
Under \eqref{30006791:sha}, Stoll's Corollary 7.4 says that
$C(\A_k)^{\operatorname{Br}}\ne\varnothing$ forces a $k$-rational divisor
class of degree one. Such a class gives a closed embedding $i:C\hookrightarrow J$,
and the same corollary gives
\[
 i\bigl(C(\A_k)^{\operatorname{Br}}_\bullet\bigr)
 =i\bigl(C(\A_k)_\bullet\bigr)\cap\overline{J(k)}.
\]
The subscript means passage to connected components at infinity.
\sref{30006791}{1} In particular, any fixed Brauer tuple can be approximated
at any finite collection of finite places by elements of $J(k)$.
None of this conclusion is being invoked without \eqref{30006791:sha}.
A rational divisor \emph{class} is sufficient for the embedding; it has
not been replaced by an assumed rational point on $C$.

\textbf{2. A height-tightness lemma.}
Fix a projective embedding $J\hookrightarrow\PP^r_k$ and its absolute
logarithmic height $h$. Enumerate the good finite places as
$v_1,v_2,\ldots$. Suppose there are $a_n\in J(k)$ whose reductions at
$v_1,\ldots,v_n$ lie on the corresponding reduction of $i(C)$ and for
which $h(a_n)\le H$, with one finite $H$ independent of $n$. Then $C(k)$
is nonempty.

To prove this, the number-field degree is fixed and Northcott finiteness
makes $\{a_n:n\ge1\}$ finite. At least one value $a$ occurs for unbounded
$n$, so it reduces onto $i(C)$ at every good finite place. If
$a\notin i(C)$, choose a homogeneous polynomial $F$ over $k$ vanishing
on $i(C)$ but with $F(a)\ne0$. Outside a finite set of places, coordinates
for $a$ and coefficients for $F$ are integral and a coordinate of $a$ is
a unit. A nonzero algebraic number $F(a)$ can vanish modulo only finitely
many such prime ideals. This contradicts all the reductions of $a$.
Hence $a=i(P)$ for a $k$-point $P$ of $C$, since $i$ is a closed
$k$-embedding.

This proof uses only ordinary height finiteness and elementary reduction;
it does not require an effective height bound, consistent with the
noneffective existence target. It also explains why compactness in the
adelic topology alone is inadequate: compactness would need to take place
inside a finite set of bounded-height rational points.

\textbf{3. A quantitative height-versus-sieve lemma over any number field.}
The following gives a stronger stopping test. Put $D=[k:\Q]$ and use
absolute values extending the standard ones on $\Q$, with weights
$n_v=[k_v:\Q_v]/D$. Thus $\sum_v n_v\log|b|_v=0$ for $b\in k^\times$.
For projective coordinates $x=(x_0,\ldots,x_r)$ set
\[
 \|x\|_v=\max_i|x_i|_v,\qquad h([x])=\sum_v n_v\log\|x\|_v.
\]
Let $F\in k[X_0,\ldots,X_r]$ be homogeneous of degree $e\ge1$.
At nonarchimedean $v$ let $c_v(F)$ be the maximum of $1$ and its coefficient
absolute values; at archimedean $v$ use the maximum of $1$ and their sum.
The finite constant
$c(F)=\sum_v n_v\log c_v(F)$ gives
$|F(x)|_v\le c_v(F)\|x\|_v^e$ everywhere.

Let $T$ be a finite set of good prime ideals where these coefficients are
integral and where the projective reduction of $[x]$ satisfies $F=0$.
Write $M_T=\prod_{\mathfrak p\in T}N\mathfrak p$. If $F(x)\ne0$, then
\begin{equation}\label{30006791:sieve}
 \log M_T\le D\bigl(eh([x])+c(F)\bigr).
\end{equation}
Indeed, after locally scaling a coordinate to a unit, divisibility by
$\mathfrak p$ gives
\[
 \frac{|F(x)|_v}{\|x\|_v^e}\le p^{-1/e_v}
 \quad(v\leftrightarrow\mathfrak p\mid p),
 \qquad
 n_v\log(p^{-1/e_v})=-D^{-1}\log N\mathfrak p.
\]
At every remaining place the same normalized expression is bounded by
$c_v(F)$. Sum its logarithm. The product formula makes the left sum
$-eh([x])$, so
$-eh([x])\le c(F)-D^{-1}\log M_T$, proving \eqref{30006791:sieve}.
This argument is invariant under rescaling $x$, and so does not assume
that the number field has principal ideal class group.

Choose finitely many homogeneous equations $F_1,\ldots,F_s$ cutting out
$i(C)$ inside the projective model of $J$, and enlarge the finite excluded
set of primes so that reduction respects these equations. Put
$e=\max_j\deg F_j$ and $c=\max_jc(F_j)$. Every
$a\in J(k)\setminus i(C)$ which survives the sieve at $T$ must satisfy
\begin{equation}\label{30006791:separation}
 h(a)\ge\frac{D^{-1}\log M_T-c}{e}.
\end{equation}
To see this, select an equation not vanishing at $a$ and apply
\eqref{30006791:sieve}, using $h(a)\ge0$. Conversely a surviving rational point
with $D(eh(a)+c)<\log M_T$ lies on $i(C)$ and yields a rational point on
$C$. This is a finite, exact certificate, not a density heuristic.

\textbf{4. The missing selection theorem.}
The attempted completion would prove the following statement for the
Jacobians and embeddings obtained in part~1:
\begin{quote}\small
For each Brauer tuple there are increasing finite sets $T_n$ of good
places and $a_n\in J(k)$ matching its reductions at $T_n$, such that
$\log M_{T_n}\to\infty$ and, for some $\eta>0$,
\[
 h(a_n)\le(1-\eta)\frac{\log M_{T_n}}{De}
 \quad\text{for arbitrarily large }n.
\]
\end{quote}
For those $n$, eventually $D(eh(a_n)+c)<\log M_{T_n}$, and the separation
lemma gives a point of $C(k)$. Bounded-height selection is a special case.
Neither the existence of these representatives nor the height saving is
proved here. Weak approximation in the \emph{closure} supplies compatible
congruences, not this archimedean control.

\textbf{5. A concrete failure of the compactness shortcut.}
In $\widehat{\Z}=\prod_p\Z_p$ take $z_2=-1$ and $z_p=0$ for every odd prime.
Every finite collection of its prime-power congruences has an integer
representative by the Chinese remainder theorem, and the congruence sets
form a compatible inverse system. There is no integer representing $z$:
an integer divisible by $3^m$ for every $m$ is zero, but zero does not
have $2$-adic coordinate $-1$.

More quantitatively, let $a_m$ simultaneously satisfy
$a_m\equiv-1\pmod{2^m}$ and $a_m\equiv0\pmod{3^m}$. Such representatives
cannot remain bounded, because any bounded subsequence repeating one
integer would contradict the preceding argument. The generic Chinese
remainder representative bound is exponential in the logarithmic modulus;
viewing these representatives as Mordell--Weil coordinates gives no
reason for their quadratic canonical heights to meet part~4. This is an
abstract model of the obstruction, not an adelic point on a particular
curve violating Brauer--Manin.

\paragraph{Stress test.}
The full Brauer group enters through Stoll's exact comparison theorem;
it has not been replaced by a finite list of Brauer classes or by local
solubility alone. The degree-one class is deduced only under the displayed
Tate--Shafarevich hypothesis. Without that hypothesis, the torsor
$\operatorname{Pic}^1(C)$ and divisible obstruction cannot be discarded.
Even granting it does \emph{not} establish the conjecture: the new
height-selection step remains independent and unproved.

The norm inequality sums over all archimedean and nonarchimedean places,
with the field-degree weights written explicitly. Thus the argument is
not restricted to $k=\Q$ or to rational points having globally primitive
integral coordinates. Finite exceptional primes and the coefficients of
the defining equations contribute the constant $c$; they are not ignored.
The strict height inequality is essential: equality in a divisibility
bound cannot force a polynomial value to be zero.

\paragraph{Outcome.}
\textbf{Outcome: Conditional reduction.}

The unconditional output is the sieve separation bound
\eqref{30006791:separation} and the bounded-height lifting lemma. Together with
both vanishing of the relevant divisible Tate--Shafarevich subgroup and
the explicit height-saving selection statement, they imply the original
Hasse-principle assertion for all number fields. Both additional global
hypotheses remain unproved. No curve is constructed with a nonempty full
Brauer set and no rational point, and no unconditional resolution is
claimed. The product-formula lemma is a notebook derivation, not a claim
of previously unknown arithmetic geometry.
% END RESEARCH ADDITION 210
\subsection{Sources}
\begin{itemize}\raggedright
\item[\textnormal{[S1]}]\hypertarget{source:30006791:1}{} Michael Stoll, Finite descent obstructions and rational points on curves, arXiv math/0606465v3.
\url{https://arxiv.org/pdf/math/0606465v3}.
{\small\textit{Catalogue formulation source.}}
% BEGIN RESEARCH REFERENCES 210
\item[\textnormal{[E1]}]\hypertarget{extra:30006791:1}{} Brendan Creutz,
\emph{Adelic Mordell--Lang and the Brauer--Manin Obstruction},
arXiv:2510.25931v2 (22 September 2026), Theorems 1.1--1.2 and Sections 4.1 and 5.
\url{https://arxiv.org/abs/2510.25931}.
{\small\textit{Consulted 27 September 2026. The unrestricted number-field assertion is not an unconditional consequence of this preprint.}}
% END RESEARCH REFERENCES 210
\end{itemize}

\end{document}
